\documentclass[11pt,a4paper]{article}

% ── Packages ──
\usepackage[utf8]{inputenc}
\usepackage[T1]{fontenc}
\usepackage{lmodern}
\usepackage[margin=1in]{geometry}
\usepackage{booktabs}
\usepackage{tabularx}
\usepackage{array}
\usepackage{graphicx}
\usepackage{xcolor}
\usepackage{enumitem}
\usepackage{amsmath,amssymb}
\usepackage{hyperref}
\usepackage{listings}
\usepackage{caption}
\usepackage{fancyhdr}
\usepackage{titlesec}

% ── Colours ──
\definecolor{headerblue}{HTML}{2C3E50}
\definecolor{accentgreen}{HTML}{27AE60}
\definecolor{warnorange}{HTML}{E67E22}
\definecolor{dangerred}{HTML}{E74C3C}
\definecolor{codebg}{HTML}{F5F5F5}
\definecolor{codeframe}{HTML}{CCCCCC}

% ── Hyperref setup ──
\hypersetup{
    colorlinks=true,
    linkcolor=headerblue,
    urlcolor=headerblue,
    citecolor=headerblue
}

% ── Section formatting ──
\titleformat{\section}
    {\Large\bfseries\color{headerblue}}{\thesection}{1em}{}
\titleformat{\subsection}
    {\large\bfseries\color{headerblue!80}}{\thesubsection}{1em}{}
\titleformat{\subsubsection}
    {\normalsize\bfseries\color{headerblue!60}}{\thesubsubsection}{1em}{}

% ── Listings style ──
\lstdefinestyle{simspec}{
    backgroundcolor=\color{codebg},
    frame=single,
    rulecolor=\color{codeframe},
    basicstyle=\ttfamily\small,
    breaklines=true,
    xleftmargin=1em,
    framexleftmargin=0.5em,
    aboveskip=1em,
    belowskip=1em,
}

% ── Header / Footer ──
\pagestyle{fancy}
\fancyhf{}
\fancyhead[L]{\small\textcolor{headerblue}{Scaling CybORG Simulation Complexity}}
\fancyhead[R]{\small\textcolor{headerblue}{\thepage}}
\renewcommand{\headrulewidth}{0.4pt}

% ── Custom column types ──
\newcolumntype{L}[1]{>{\raggedright\arraybackslash}p{#1}}
\newcolumntype{C}[1]{>{\centering\arraybackslash}p{#1}}

% ══════════════════════════════════════════════════════════════
\begin{document}

% ── Title ──
\begin{center}
    {\LARGE\bfseries\color{headerblue} Scaling CybORG Simulation Complexity:\\[4pt]
    From Scenario\,1b to Enterprise-Scale Cyber Defence}\\[12pt]
    {\large Vasanth}\\[4pt]
    {\normalsize February 2026}
\end{center}

\vspace{1em}
\hrule height 0.5pt
\vspace{1.5em}

% ══════════════════════════════════════════════════════════════
\section{Introduction}

This document analyses the complexity of the current CybORG Scenario\,1b simulation
used for training reinforcement learning (RL) agents for autonomous cyber defence,
and describes what a simulation \textbf{one order of magnitude higher} in complexity
would look like across six key axes: network scale, attacker sophistication,
defender action space, observation realism, reward structure, and multi-agent
organisation.

% ══════════════════════════════════════════════════════════════
\section{Current vs.\ One Order Higher}

% ── 2.1 Network Scale ──
\subsection{Network Scale}

\begin{table}[h!]
\centering
\caption{Network scale comparison}
\label{tab:network}
\begin{tabularx}{\textwidth}{@{} l X X @{}}
\toprule
\textbf{Aspect} & \textbf{Current (Scenario\,1b)} & \textbf{$\mathbf{10\times}$ Higher} \\
\midrule
Hosts & 13 & 100--200 (departmental segments, DMZs, cloud VMs, IoT endpoints) \\
Subnets & 3 (User, Enterprise, Operational) & 15--30 (HR, Finance, R\&D, OT floor, SCADA, cloud VPC, VPN gateway, partner extranet, IoT) \\
Inter-subnet routing & Flat ACLs (\texttt{in: all/None}) & Dynamic firewall rules, micro-segmentation, NAT, VPN tunnels, SDN policies \\
Topology & Static, fully known & Partially observable; hosts spin up/down (auto-scaling), DHCP-assigned IPs \\
\bottomrule
\end{tabularx}
\end{table}

\noindent
\textbf{Observation space impact:} $\sim$11,293 $\;\longrightarrow\;$ $\sim$100K--500K
(or a graph neural network representation would replace the flat vector).

% ── 2.2 Attacker Sophistication ──
\subsection{Attacker Sophistication}

\begin{table}[h!]
\centering
\caption{Attacker sophistication comparison}
\label{tab:attacker}
\begin{tabularx}{\textwidth}{@{} l X X @{}}
\toprule
\textbf{Aspect} & \textbf{Current} & \textbf{$\mathbf{10\times}$ Higher} \\
\midrule
Red agents & 2 (B\_line, Meander) & 8--12 distinct adversary profiles \\[3pt]
Attack techniques & 6 actions (discover, exploit, escalate, impact) &
    30--50 MITRE ATT\&CK techniques: spear-phishing, credential stuffing,
    LOLBins, supply-chain implants, fileless malware, DNS tunnelling, DGA C2 \\[3pt]
Multi-stage campaigns & Single kill chain &
    Branching kill chains with fallback strategies \\[3pt]
Stealth / evasion & None---Red is oblivious to Blue &
    Adaptive adversary: detects Blue's monitoring, slows down, uses encrypted
    channels, timestamps attacks during business hours to blend with Green \\[3pt]
Persistence & Removed by \texttt{Restore} &
    Multiple persistence mechanisms per host (registry keys, scheduled tasks,
    bootkits)---partial removal possible \\[3pt]
Lateral movement & Linear (User $\to$ Enterprise $\to$ Operational) &
    Graph-based: any compromised host can pivot to any reachable host;
    credential reuse, pass-the-hash, Kerberoasting \\[3pt]
Multi-agent Red & 1 attacker &
    2--3 coordinated Red teams (e.g.\ initial access broker hands off
    to ransomware operator) \\
\bottomrule
\end{tabularx}
\end{table}

% ── 2.3 Defender Action Space ──
\subsection{Defender Action Space}

\begin{table}[h!]
\centering
\caption{Defender action space comparison}
\label{tab:defender}
\begin{tabularx}{\textwidth}{@{} l X X @{}}
\toprule
\textbf{Aspect} & \textbf{Current} & \textbf{$\mathbf{10\times}$ Higher} \\
\midrule
Blue actions & 6 types (Monitor, Analyse, Remove, Misinform, Restore, Sleep) &
    20--30 action types \\[3pt]
New actions & --- &
    \texttt{IsolateHost}, \texttt{BlockIP}, \texttt{UpdateFirewallRule},
    \texttt{DeployHoneypot}, \texttt{RotateCredentials},
    \texttt{PatchVulnerability}, \texttt{EnableMFA},
    \texttt{TakeForensicImage}, \texttt{EscalateToHuman},
    \texttt{DeployEDR}, \texttt{SinkholeDNS}, \texttt{RevokeToken},
    \texttt{SegmentSubnet} \\[3pt]
Action targeting & 1 host at a time &
    Batch actions (e.g.\ patch all hosts in a subnet), conditional actions
    (if $X$ then $Y$), resource budgets (max 3 actions/turn) \\[3pt]
Action duration & Instant (1 step) &
    Multi-step: forensic image $=$ 3 turns, patching $=$ 2 turns with
    downtime window \\[3pt]
Action cost model & Fixed (\texttt{Restore} $= -1$) &
    Variable costs: monetary budget, analyst time, user-disruption SLA,
    downtime windows \\
\bottomrule
\end{tabularx}
\end{table}

\noindent
\textbf{Action space impact:}
$\text{Discrete}(54) \;\longrightarrow\; \text{Discrete}(300\text{--}600)$
or continuous / multi-discrete.

% ── 2.4 Observation Realism ──
\subsection{Observation Realism}

\begin{table}[h!]
\centering
\caption{Observation realism comparison}
\label{tab:obs}
\begin{tabularx}{\textwidth}{@{} l X X @{}}
\toprule
\textbf{Aspect} & \textbf{Current} & \textbf{$\mathbf{10\times}$ Higher} \\
\midrule
Sensor feeds & Single flat vector from \texttt{FixedFlatWrapper} &
    Multiple heterogeneous streams: SIEM alerts (structured), NetFlow
    (time-series), EDR telemetry (event logs), DNS query logs,
    authentication logs \\[3pt]
Noise model & Green agent does ping / port scan &
    Realistic traffic distributions: 10,000+ benign events/turn with rare
    anomalies; false-positive rates of 1--5\%; alert-fatigue modelling \\[3pt]
Observation delay & Instant (same turn) &
    Delayed / lossy: network tap misses 2\% of packets, SIEM has 30\,s
    ingestion lag, some hosts report only every $N$ turns \\[3pt]
Deception by Red & None &
    Red plants false indicators: fake log entries, timestomping,
    anti-forensics that corrupt Blue's \texttt{Analyse} results \\
\bottomrule
\end{tabularx}
\end{table}

% ── 2.5 Reward and Objective Complexity ──
\subsection{Reward and Objective Complexity}

\begin{table}[h!]
\centering
\caption{Reward and objective complexity comparison}
\label{tab:reward}
\begin{tabularx}{\textwidth}{@{} l X X @{}}
\toprule
\textbf{Aspect} & \textbf{Current} & \textbf{$\mathbf{10\times}$ Higher} \\
\midrule
Reward signal & Per-host penalty based on Red access &
    Multi-objective: confidentiality loss (exfiltration volume), availability
    (SLA uptime~\%), integrity (tampered records), compliance (audit score),
    reputation damage \\[3pt]
Business context & Hosts equally interchangeable within tier &
    Asset criticality graph: ERP database $= 100\times$ a user laptop;
    CEO's laptop has board-level secrets; OT controller failure $\Rightarrow$
    physical safety risk \\[3pt]
Time pressure & Fixed episode (30 / 50 / 100 steps) &
    Open-ended episodes with escalating threat level, shift changes
    (night shift $=$ fewer analysts), maintenance windows \\[3pt]
Opponent co-adaptation & Red is fixed &
    Red adapts to Blue's strategy over episodes (game-theoretic,
    non-stationary) \\
\bottomrule
\end{tabularx}
\end{table}

% ── 2.6 Multi-Agent and Organisational ──
\subsection{Multi-Agent and Organisational}

\begin{table}[h!]
\centering
\caption{Multi-agent and organisational comparison}
\label{tab:multiagent}
\begin{tabularx}{\textwidth}{@{} l X X @{}}
\toprule
\textbf{Aspect} & \textbf{Current} & \textbf{$\mathbf{10\times}$ Higher} \\
\midrule
Blue agents & 1 autonomous agent &
    SOC team of 3--5 agents with roles (Tier\,1 triage, Tier\,2
    investigation, Tier\,3 incident response); communication overhead;
    coordination \\[3pt]
Human-in-the-loop & None &
    Escalation to human analyst with response delay; approval required
    for high-impact actions (isolating a production server) \\[3pt]
Cross-organisation & Single network &
    Federated defence: share threat intel with partner organisations,
    coordinated response across supply chain \\
\bottomrule
\end{tabularx}
\end{table}

% ══════════════════════════════════════════════════════════════
\section{Concrete Specification}

The following listing summarises what a $10\times$ simulation would look like
as a concrete scenario specification:

\begin{lstlisting}[style=simspec]
Simulation: "Enterprise SOC Scenario"
--------------------------------------
Network:      150 hosts, 20 subnets, 3 cloud VPCs, 1 OT segment
Obs space:    ~250,000-d (or graph: 150 nodes x 200 features + edge matrix)
Action space: MultiDiscrete([25, 150, 5])  ->  action_type x target x priority
Red agents:   8 profiles (APT28-style, ransomware, insider threat, ...)
Blue agents:  3 cooperating agents with role specialisation
Episode:      500-1000 steps (simulating ~1 week of real time)
Sensors:      SIEM (delayed), EDR (per-host), NetFlow (sampled), DNS logs
Budget:       100 analyst-hours/episode, $50K incident response budget
Objectives:   Minimise data loss, maintain 99.5% uptime SLA,
              comply with audit requirements, minimise response cost
\end{lstlisting}

% ══════════════════════════════════════════════════════════════
\section{Key Technical Challenges at This Scale}

\begin{table}[h!]
\centering
\caption{Technical challenges introduced by $10\times$ scaling}
\label{tab:challenges}
\begin{tabularx}{\textwidth}{@{} l X @{}}
\toprule
\textbf{Challenge} & \textbf{Why It Is Hard} \\
\midrule
Observation representation &
    Flat vectors do not scale to 150+ hosts. Requires \textbf{graph neural
    networks} (GNN) or \textbf{attention-based architectures} to handle
    variable-size, relational observations. \\[4pt]
Action space explosion &
    $25 \times 150 = 3{,}750$ actions. Requires \textbf{hierarchical RL},
    \textbf{action abstraction}, or \textbf{parameterised action spaces}. \\[4pt]
Multi-agent coordination &
    3--5 Blue agents must share information and avoid redundant actions
    $\Rightarrow$ \textbf{MARL} (QMIX, MAPPO, communication protocols). \\[4pt]
Non-stationary opponent &
    Adaptive Red requires \textbf{meta-learning}, \textbf{robust RL}, or
    \textbf{game-theoretic approaches} (PSRO, fictitious play). \\[4pt]
Sample efficiency &
    Complex simulation is slow to run. Requires \textbf{model-based RL},
    \textbf{transfer learning} from simpler scenarios, or \textbf{offline RL}
    from real SOC data. \\[4pt]
Sim-to-real gap &
    More complex simulation $\neq$ more realistic. Requires
    \textbf{domain randomisation} and \textbf{calibration against real
    network telemetry}. \\
\bottomrule
\end{tabularx}
\end{table}

% ══════════════════════════════════════════════════════════════
\section{Recommended Progression}

Rather than jumping directly to $10\times$ complexity, a practical staged
progression from the current Scenario\,1b setup is recommended:

\begin{enumerate}[leftmargin=2em, itemsep=6pt]
    \item \textbf{$\mathbf{3\times}$ complexity} ---
        Add 5 more hosts, 2 more subnets, 3 Red agent profiles.
        Introduce \texttt{IsolateHost} and \texttt{PatchVulnerability} actions.
        Add observation delay (1-step lag).

    \item \textbf{$\mathbf{5\times}$ complexity} ---
        Switch to \textbf{graph-based observation} (adjacency matrix $+$
        node features).  Add multi-step actions.  Introduce analyst time
        budget.  Add adaptive Red agent.

    \item \textbf{$\mathbf{10\times}$ complexity} ---
        Multi-agent Blue team.  Full MITRE ATT\&CK action set.  Realistic
        traffic noise model.  Multi-objective reward.
\end{enumerate}

\vspace{0.5em}
\noindent
The current DQN result demonstrates that the core ``don't collapse to one
action'' problem has been solved.  The next meaningful step is
\textbf{moving from flat observations to graph-structured representations}
--- that single change would unlock scalability to much larger networks
without the 11,293-dimensional sparse vector problem.

% ══════════════════════════════════════════════════════════════
\vfill
\hrule height 0.5pt
\vspace{0.5em}
\begin{center}
\small\textcolor{gray}{Generated for the RL\_CYB project -- February 2026}
\end{center}

\end{document}
